% Sources breakdown bar chart
\begin{figure}[htbp]
\centering
\begin{tikzpicture}
\begin{axis}[
  xbar,
  width=10cm, height=7cm,
  bar width=12pt,
  xlabel={Number of HTML pages / sources},
  ytick={1,2,3,4,5,6},
  yticklabels={
    Powder Cryst.\ book (opt.),
    Programmer's Guide,
    Help manual,
    Home / Install / Dev,
    Tutorials (dynamic),
    --- total default ---
  },
  yticklabel style={font=\small},
  xmin=0, xmax=215,
  xtick={0,25,50,75,100,125,150,175,200},
  nodes near coords,
  nodes near coords align={horizontal},
  every node near coord/.style={font=\small},
  axis lines*=left,
  tick align=outside,
  enlarge y limits=0.15,
]
\addplot[fill=argblue!70, draw=argblue]
  coordinates {(65,5) (22,4) (42,3) (23,2) (179,1)};
\addplot[fill=apsorange!70, draw=apsorange,
  nodes near coords style={xshift=2pt}]
  coordinates {(152,6)};
\end{axis}
\end{tikzpicture}
\caption{Documentation sources indexed by \textit{Query-GSAS} by category.
  Blue bars show the default index (152 HTML sources).
  The \textit{Powder Diffraction Crystallography} textbook (179 HTML pages,
  orange) is included via \texttt{--book} and not indexed by default.
  Tutorials are resolved dynamically from the GSAS-II repository at
  ingest time.}
\label{fig:sources}
\end{figure}